

\documentclass[11pt,english,ngerman, headsepline]{scrreprt}
\usepackage{lmodern}
\renewcommand{\sfdefault}{lmss}
\renewcommand{\ttdefault}{lmtt}
\usepackage[T1]{fontenc}

\usepackage{listings}

\usepackage[utf8]{inputenc} \usepackage[a4paper]{geometry}
\geometry{verbose,tmargin=3cm,bmargin=3cm,lmargin=3cm,rmargin=2.75cm,headheight=1cm,headsep=0.666cm,footskip=1cm}
\setcounter{secnumdepth}{3}
\setcounter{tocdepth}{3}
\setlength{\parskip}{\medskipamount}
\setlength{\parindent}{0pt}
\usepackage{babel}

\usepackage{verbatim} 
\usepackage{float}  
\usepackage{url}
\usepackage{graphicx}
\usepackage{setspace}
\usepackage[square,sort,numbers]{natbib}
\usepackage[utf8]{inputenc}
\usepackage{graphicx}
\usepackage[xindy,toc]{glossaries}



\setstretch{1.4}
\usepackage[usenames,dvipsnames]{color}
\usepackage[unicode=true, 
 bookmarks=true,bookmarksnumbered=false,bookmarksopen=true,bookmarksopenlevel=2,
 breaklinks=false,pdfborder={0 0 0},backref=false,colorlinks=false]
 {hyperref}
\hypersetup{pdftitle={Vergleich und Evaluation zwischen modernen und traditionellen Datenbankkonzepten 
unter den Gesichtspunkten Skalierung, Abfragem�glichkeit und Konsistenz.},
 pdfauthor={Nils M. Petersohn}}
 
\makeatletter

%%%%%%%%%%%%%%%%%%%%%%%%%%%%%% LyX specific LaTeX commands.
\providecommand{\LyX}{L\kern-.1667em\lower.25em\hbox{Y}\kern-.125emX\@}
%% Because html converters don't know tabularnewline
\providecommand{\tabularnewline}{\\}

%%%%%%%%%%%%%%%%%%%%%%%%%%%%%% Textclass specific LaTeX commands.
\newenvironment{lyxcode}
{\par\begin{list}{}{
\setlength{\rightmargin}{\leftmargin}
\setlength{\listparindent}{0pt}% needed for AMS classes
\raggedright
\setlength{\itemsep}{0pt}
\setlength{\parsep}{0pt}
\normalfont\ttfamily}%
 \item[]}
{\end{list}}

%%%%%%%%%%%%%%%%%%%%%%%%%%%%%% User specified LaTeX commands.
%% Flexibles Seitenlayout
\usepackage[automark]{scrpage2}

%% Mehrspaltenlayout ermöglichen
\usepackage{multicol}

%% Unterst\"utzung f\"ur Farben
\usepackage{color}

%% Schönere Tabellen
\usepackage{booktabs, longtable}

%% Schönerer Blocksatz
\usepackage{microtype}
%% Mehr Platz zwischen Überschrift und Tabelle
\newcommand{\@ldtable}{}
\let\@ldtable\table
\renewcommand{\table}{ %
    \setlength{\@tempdima}{\abovecaptionskip} %
    \setlength{\abovecaptionskip}{\belowcaptionskip} %
    \setlength{\belowcaptionskip}{\@tempdima} %
    \@ldtable %
}

%% Verschiedene Symbole und Zeichen wie (c), ¤
\usepackage{textcomp}

%% Fehlerkorrektur f\"ur Marginalien
\usepackage{fixltx2e }%,mparhack

%% Deutsche Kurzfassung und englisches Abstract auf eine Seite
\renewenvironment{abstract}{
    \@beginparpenalty\@lowpenalty
        \begin{center}
            \normalfont\sectfont\nobreak\abstractname
        \end{center}
    \@endparpenalty\@M
}{
    \par
}

%% Alle Seiten vor dem Inhaltsverzeichnis sind römisch nummeriert
\pagenumbering{roman}
\let\myTOC\tableofcontents
\renewcommand\tableofcontents{
    \begin{spacing}{1.1}
    \myTOC
    \end{spacing}
    \clearpage
    \pagenumbering{arabic}
}

%% Kopfzeile um Logo erg\"anzen
\clearscrheadfoot
\ohead{\\\headmark}
\ihead{\includegraphics[scale=0.25]{pics/FH-Logo-only-gray}}%\pagemark}
\ofoot[\pagemark]{\pagemark}

%% Randnotizen anpassen
\setlength{\marginparwidth}{22mm}
\let \oldmarginpar = \marginpar
\renewcommand{\marginpar}[1]{%
    \-\oldmarginpar[\raggedleft\footnotesize\sf #1]%
        {\raggedright\footnotesize\sf #1%
    }}

%% Zitate am Kapitelanfang
%\usepackage{epigraph}
%\setlength{\epigraphwidth}{9cm}

%% Code-Block-Formatierung
\lstdefinestyle{default}{ %
    backgroundcolor={\color[rgb]{0.95,0.95,0.95}}, %
    basicstyle={\small\ttfamily}, %
    breaklines=true, %
    frame=l, %
    language={[Sharp]C}, %
    lineskip=-0.1pt, %
    numbers=left, %
    rulecolor={\color[rgb]{0.5,0.5,0.5}}, %
    xleftmargin={0.75cm}, %
    xrightmargin={0cm} %
}
\lstdefinelanguage{JavaScript} {
	morekeywords={
		break,const,continue,delete,do,while,export,for,in,function,
		if,else,import,in,instanceOf,label,let,new,return,switch,this,
		throw,try,catch,typeof,var,void,with,yield
	},
	sensitive=false,
	morecomment=[l]{//},
	morecomment=[s]{/*}{*/},
	morestring=[b]",
	morestring=[d]'
}

\lstset{
	frame=tb,
	framesep=5pt,
	basicstyle=\footnotesize\ttfamily,
	showstringspaces=false,
	keywordstyle=\ttfamily\bfseries\color{CadetBlue},
	identifierstyle=\ttfamily,
	stringstyle=\ttfamily\color{OliveGreen},
	commentstyle=\color{GrayBlue},
	rulecolor=\color{Gray},
	xleftmargin=5pt,
	xrightmargin=5pt,
	aboveskip=\bigskipamount,
	belowskip=\bigskipamount
}
\makeatother

\setlength{\parindent}{5pt} 
\setlength{\parskip}{1ex}

\parindent 0pt


\begin{document}


\titlepage

\begin{center}
\includegraphics[width=5cm]{pics/FH-Logo}\vspace{0.5cm}

\par\end{center}

\noindent \begin{center}
\textsf{\textbf{\Large Fachbereich Informatik und Medien}}\\
\textsf{\large }\\
\vspace{1cm}

\par\end{center}

\noindent \begin{center}
\textsf{\textbf{\huge Apache Camel  \\ Anwendung - Erdbebendaten}}\textsf{}\\
\textsf{}\\
\textsf{\Large Systemintegration 1. Semester Master Informatik \\ Prof.
Dr. Preuss }


\par\end{center}{\Large \par}

\vspace{2cm}


\noindent \begin{center}
{\huge }\begin{tabular}{rl}
Vorgelegt von: & Robert Stümer, Nils Petersohn, Steffen
Reinkensmeier\tabularnewline am: & 27.01.2011.\tabularnewline
\end{tabular}
\par\end{center}{\huge \par}

\vspace{1cm}




\noindent \begin{center}
\medskip{}
\begin{tabular}{rl}
\tabularnewline
\end{tabular}
\par\end{center}

\newpage{}

%
\begin{comment}
leere Seite nach dem Titelblatt

dann Aufgabenstellung
\end{comment}


\begin{comment}
%\pagestyle{scrheadings}    %Kopfzeile ein

\noindent \begin{center}
\textsf{\textbf{\large Selbstst\"andigkeitserkl\"arung}}
\par\end{center}{\large \par}

\noindent Hiermit erkl\"are ich, dass ich die vorliegende Arbeit zum
Thema

\smallskip{}


\noindent \begin{center}
\textsf{Medienverwaltung mit Eclipse und Equinox}
\par\end{center}

\smallskip{}


\noindent vollkommen selbstst\"andig verfasst und keine anderen als
die angegebenen Quellen und Hilfsmittel benutzt sowie Zitate kenntlich
gemacht habe. Die Arbeit wurde in dieser oder \"ahnlicher Form noch
keiner anderen Pr\"ufungsbeh\"orde vorgelegt.

\medskip{}


\noindent Brandenburg/Havel, den dd.MM.yyyy

\vspace{1.7cm}


\noindent Unterschrift 

\selectlanguage{ngerman}
\end{comment}
%\newpage{}

%\noindent \begin{center}
%\textsf{\textbf{\large Danksagung}}
%\par\end{center}{\large \par}


%\newpage{}
%\begin{abstract}

%\end{abstract}

%\selectlanguage{english}%
%\begin{abstract}
%This is a second abstract in english.

%\newpage{}
%\end{abstract}

\selectlanguage{ngerman}%
\tableofcontents{}

\pagestyle{scrheadings}    %Kopfzeile ein


%==========================================================================
% DOCUMENT START 
%==========================================================================

\chapter{Introduction / Motivation / Project}
 
 Systemintegration is the part of a software architectur. It helps to connect
 components or subsystems together. Certain patterns are used in the industry
 today. Using and learning EIP (``Enterprise Integration Pattterns'') with
 Apache Camel is the goal of this Project.\\
 
 The Example for this project is all about earthquake data from around the
 world. The Application is able to read earthquake data from various rss Feeds
 and processes it. During the processing the data will be in form of XML and
 Java Objects. The data will be enriched, splitted, sorted, aggregated,
 normalized, marshalled umarshalled and finally provided again in form of a
 restful service.\\
 
 The specified task is as follows:
 
 \begin{enumerate}
   
 
 \item Read Earthquake Data continously from those two RSS Streams
 \begin{itemize}
  	\item \url{http://geofon.gfz-potsdam.de/db/eqinfo.php?fmt=rss}
  	\item \url{http://earthquake.usgs.gov/eqcenter/catalogs/eqs1day-M2.5.xml}
  \end{itemize}
  \item enrich this data with other related information like the weather in this
  area at this time. Data can be from here: \url{http://www.programmableweb.com}.
  \item sort the earthquakes by the earthparts where they appear
  \item if the earthquake has a strength of more than ``M 5.5'' than send an
  formated warning email to an email adress.
  \item  provide this data via a Restful interface in form of a list of the
  earthparts with an xlink to detailed information of the earthquakes.
  
 \end{enumerate}

\chapter{Information Collection and Normalization}\label{chapter:normalization}

The Normalizer \cite{hohpe2003enterprise} (Pic. \ref{normalizerPic}) routes each
message type through a custom Message Translator so that the resulting messages match a common format. 
The different message formates are collected in one channel
(``$direct:collectorChannel$'') and than routed according to their xml format.
The Routing is accomlished with a xpath condition which routes the messages to
the specified Message Translators \cite{hohpe2003enterprise} (Pic.
\ref{translatorPic}). Here The Translator is implemented with a XSLT formatter
\url{http://www.w3.org/TR/xslt}.
The relative source code is listed in \ref{normalizerTranslator.java}


  \begin{figure}[h!]
	\begin{center}
	\includegraphics[width=0.9\textwidth]{pics/NormalizerDetail.jpg}
	\end{center}
	\caption{Normalizer Pattern \cite{hohpe2003enterprise}}
	\label{normalizerPic} 
   \end{figure}
   
   
   \begin{figure}[h!]
	\begin{center}
	\includegraphics[width=0.9\textwidth]{pics/MessageTranslator.jpg}
	\end{center}
	\caption{Translator Pattern \cite{hohpe2003enterprise}}
	\label{translatorPic} 
   \end{figure}



\lstinputlisting[caption={Normalizer and Translator Source Code},
label={normalizerTranslator.java},style=default]{code/normalizerTranslator.java}




\chapter{Aggregating Information Sources} 



\begin{quote}
``The Aggregator is a special filter that recieves a strem of messages an
didentifies messages that are correlated. Once a complete set of messages has
been recied, the Aggregator collects information from each correlated message
and publishes a single, aggregated message to the output channel for further
processing.'' \cite{hohpe2003enterprise}
\end{quote}

  \begin{figure}[h!]
	\begin{center}
	\includegraphics[width=0.9\textwidth]{pics/Aggregator.jpg}
	\end{center}
	\caption{Aggregator Pattern \cite{hohpe2003enterprise}}
	\label{utilityTree} 
   \end{figure}


Camel provides a method named $aggregate$ with two parameters, first the
identifier (the message header which was defined before in chapter
\ref{chapter:normalization}) of the messages to aggregate  and second the
strategy (as a Object: $SpecialAggregationStrategy$ which implements\\
$org.apache.camel.processor.aggregate.AggregationStrategy$). Code listing
\ref{aggregateinputstreams.java} shows the implementation in this project for
the normalized messages from Chapter \ref{chapter:normalization}. Code listing
\ref{aggregationstrategy.java} shows the strategy which is fairly simple.


\lstinputlisting[caption={Aggregator},
label={aggregateinputstreams.java},style=default]{code/aggregateinputstreams.java}


\lstinputlisting[caption={Aggregation Strategy},
label={aggregationstrategy.java},style=default]{code/aggregationstrategy.java}


\chapter{Enrich Data with related Information}



\chapter{Sortierung und speicherung der Daten}
\chapter{Emailbenachrichtigung}
\chapter{REST Schnittstelle zu den Daten}
\section{Liste der Erdbeben je Erdteil mit xlink}
\section{Detailinformationen zu einem Erdbeben}

%\chapter{verwendete EAI-Pattern}



 
%==========================================================================
% DOCUMENT END 
%==========================================================================

\bibliographystyle{alphadin}
\clearpage\addcontentsline{toc}{chapter}{\bibname}\bibliography{jabref}


\appendix
\renewcommand{\theequation}{A-\arabic{equation}}

\setcounter{equation}{0}  % reset counter \chapter*{Anhang}  % use *-form to




\listoffigures
\end{document}
